\subsection{Finite Element Methods}\label{subseq:finite_element_methods}

In this subsection we will be looking at the discretization techniques Finite Element Methods entail, to tackle spatial discretization of functions of the form as seen in Equation~\ref{eq:general_pde_numerical_methods}. The core idea of a finite element method is to approximate the solution by a sum of finite element base functions, which are piecewise polynomial functions defined on a mesh of the domain. The finite element method then seeks to find the coefficients of these base functions that best approximate the solution to the PDE.

Say we have a PDE of the form:
\begin{equation}\label{eq:num_general_pde_numerical_methods}
    \partial_t\rho = -v\partial_x\rho \quad \text{in } \Omega
\end{equation}
where $\rho$ is the unknown function, $v$ is the flow speed, and $\Omega$ is the domain of interest. The finite element method aims to find the closest approximation of $\rho$ by a linear combination of finite element base functions $\psi^i$:
\begin{equation}
    \rho \approx \sum_i \rho^i \psi^i
\end{equation}
where $\rho^i$ are the coefficients to be determined. The finite element base functions $\psi^i$ are typically chosen to be piecewise polynomials that are defined on a mesh of the domain $\Omega$. The choice of the base functions and the mesh will depend on the specific problem and the desired accuracy of the solution. Now the way we will find the coefficients $\rho^i$ is by using test functions and testing whether they satisfy the PDE in a weak sense, so:

\begin{equation}
    \int_\Omega w^j \left( \partial_t\rho + v\partial_x\rho \right) dx = 0 \quad \forall j
\end{equation}
for a set of test functions $w^j$. In the Galerkin method, we choose the test functions to be the same as the finite element base functions, so $w^j = \psi^j$. 
\cite{toshniwal2024afe}

\subsubsection*{Lagrange Polynomials}